\chapitrepage{Réalisation du système}

\begin{objectifbox}
Ce chapitre présente l'environnement de développement, la réalisation de la
chaîne de traitement, les interfaces principales, le déploiement, la stratégie de
test et les résultats observés sur la version synchronisée du projet.
\end{objectifbox}

\section{Introduction}

La réalisation a été conduite par incréments. Nous avons d'abord sécurisé la
lecture et la consolidation des ventes, puis construit les prévisions mensuelles
et journalières. Les règles de fêtes, de commandes et de nomenclatures ont ensuite
été intégrées avant la création du dashboard et de l'automatisation. Chaque étape
a produit des fichiers contrôlables indépendamment de l'interface.

\section{Environnement de développement}

\subsection{Environnement matériel}

Le développement et la validation finale ont été réalisés sur la configuration du
tableau~\ref{tab:env-materiel}. Le calcul reste principalement contraint par le
nombre de produits et les backtests répétés ; il ne nécessite pas de carte
graphique dédiée.

\begin{table}[H]
\centering
\caption{Configuration matérielle de développement et de validation}
\label{tab:env-materiel}
\begin{tabularx}{0.86\textwidth}{>{\bfseries}p{4.2cm}Y}
\toprule
Composant & Configuration \\
\midrule
Processeur & AMD Ryzen 5 7500F, 6 cœurs \\
Mémoire vive & 15,6 Go \\
Stockage principal & SSD, capacité d'environ 930 Go \\
Architecture & 64 bits \\
Système d'exploitation & Windows 11 Professionnel \\
\bottomrule
\end{tabularx}
\end{table}

Le serveur cible documenté est un environnement Windows relié à SQL Server
Express 2014 et au progiciel Elyx Resto 11.15.5. La connexion utilise
l'authentification Windows ; aucun mot de passe n'est inscrit dans le code.

\subsection{Environnement logiciel}

\begin{table}[H]
\centering
\caption{Environnement logiciel utilisé lors de la validation}
\label{tab:env-logiciel}
\begin{tabularx}{\textwidth}{p{3.4cm}p{2.4cm}Y}
\toprule
\textbf{Outil} & \textbf{Version} & \textbf{Usage} \\
\midrule
Python & 3.11.9 & Langage principal et orchestration. \\
pandas & 2.2.2 & Lecture, nettoyage, agrégation et transformation tabulaire. \\
NumPy & 1.26.4 & Calcul numérique et manipulation des vecteurs de prévision. \\
statsmodels & 0.14.6 & Holt--Winters, Theta et ARIMA. \\
scikit-learn & 1.8.0 & Modèle global de gradient boosting. \\
Dash & 4.3.0 & Application web et callbacks interactifs. \\
Plotly & 6.7.0 & Graphiques interactifs. \\
openpyxl & 3.1.5 & Lecture et génération de classeurs Excel. \\
PyODBC & dépendance requise & Connexion au SQL Server Elyx. \\
Pytest & dépendance de test & Exécution de la suite automatisée. \\
Git / GitHub & dépôt distant & Versionnement et déploiement du code. \\
\bottomrule
\end{tabularx}
\end{table}

\section{Réalisation de la chaîne de données}

\subsection{Import historique}

Le script de conversion reconstruit le fichier
\texttt{ventes\_journalieres.csv} à partir des exports bruts. Il traite l'encodage
hérité, associe les familles, retire les sous-totaux et agrège les lignes par date
et produit. Les colonnes finales sont volontairement stables : Date, Code,
Produit, Famille, Quantité et CA TTC. Cette stabilité permet aux moteurs et au
dashboard de partager la même source de vérité journalière.

\subsection{Export SQL défensif}

Le module \texttt{exporter\_ventes\_sql.py} recherche les tables
\texttt{IMPUTATION\_<site>}, sélectionne les journées clôturées et agrège les
lignes de ticket. Les ventes SQL remplacent les lignes de même date dans le CSV,
mais laissent les autres dates intactes. Avant toute écriture, le module vérifie
la présence des colonnes attendues, le volume retourné et la cohérence des totaux.
Si les caisses ne sont pas encore raccordées ou si la requête ne renvoie rien,
l'opération devient un no-op : le fichier historique n'est pas écrasé.

\subsection{Configuration externalisée}

Les éléments métier sont stockés dans \texttt{data/*.json}. Cette décision évite
de mélanger les paramètres et le code. La version auditée contient notamment 40
périodes de fêtes marocaines, 848 choix de modèle, 270 entrées de recettes dans le
référentiel principal et 86 groupes d'alias matières. Le dashboard reconnaît 261
recettes comme exactes et correctement rattachées parmi 1\,109 produits prévus ;
les autres sont signalées comme estimées ou manquantes.

\section{Réalisation des prévisions}

\subsection{Pipeline mensuel}

Le point d'entrée \texttt{main.py} appelle \texttt{pipeline.run()}. Le pipeline
crée un dossier d'export daté, charge et nettoie les données, calcule les
prévisions agrégées, boucle sur les produits, applique le modèle global si
nécessaire, réconcilie les niveaux, ajoute les ajustements métier, estime
l'incertitude et calcule les besoins matières.

Le pseudo-code suivant résume l'enchaînement sans exposer les détails de chaque
modèle :

\begin{lstlisting}[style=paulpython,caption={Pseudo-code du pipeline mensuel},label={lst:pipeline}]
donnees = charger_et_nettoyer_sources()
historique = agreger_par_mois_et_par_produit(donnees)

for produit in produits:
    modele = modele_selectionne(produit)
    prevision[produit] = prevoir(modele, historique[produit])

prevision = appliquer_modele_global_si_retenu(prevision)
prevision = reconcilier_sur_agregat(prevision)
prevision = appliquer_fetes_evenements_commandes(prevision)
plan = ajouter_intervalles_et_marge_de_securite(prevision)
besoins = exploser_nomenclatures(plan)
exporter(plan, besoins, validations, graphiques)
\end{lstlisting}

\subsection{Prévision journalière et suivi}

Le module journalier génère un enregistrement pour chaque date et produit sur 62
jours. Il conserve la prévision centrale, la quantité recommandée, la fiabilité et
la part provenant d'une commande client. Le module \texttt{suivi.py} reconstruit
ensuite ce que le système aurait prévu au cours des 28 derniers jours en
n'utilisant, à chaque origine, que les observations disponibles auparavant. Cette
méthode empêche toute fuite d'information future.

\subsection{Gestion de l'incertitude}

La recommandation ne se limite pas à la prévision moyenne. Pour un niveau de
service cible de 95\,\%, un quantile de la loi normale transforme l'écart-type de
l'erreur en marge. Lorsque l'historique est insuffisant, une marge prudente par
défaut est appliquée et l'étiquette de fiabilité reste explicite. La quantité
recommandée répond ainsi à un objectif anti-rupture, tandis que la prévision
centrale reste disponible pour l'analyse.

\section{Réalisation du calcul matières}

\subsection{Recettes et produits semi-finis}

Les recettes sont représentées par des dictionnaires produit--ingrédient. Une
recette peut référencer un produit semi-fini, par exemple une pâte, une crème ou
une sauce. La fonction d'explosion parcourt ces références récursivement, protège
contre les cycles et additionne les matières identiques après résolution des
alias. Les derniers changements du dépôt renforcent cette logique pour les
préparations maison de cuisine et les libellés contenant un code ou une marque
interne.

Le workflow de validation par le chef comprend deux scripts. Le premier génère un
classeur par catégorie avec la recette disponible et sa provenance ; le second
importe uniquement les lignes explicitement validées, sauvegarde l'ancien JSON et
met à jour la provenance. Cette séparation réduit le risque qu'une estimation soit
promue automatiquement en donnée exacte.

\subsection{Exports opérationnels}

Le pipeline produit notamment :

\begin{itemize}
  \item \texttt{plan\_production\_securise.csv}, avec prévision, fourchette,
  recommandation, fiabilité et modèle ;
  \item \texttt{besoins\_ingredients\_planifies.csv}, avec date, ingrédient et
  quantité requise ;
  \item \texttt{previsions\_PAUL\_2026.xlsx}, classeur de synthèse ;
  \item \texttt{bon\_de\_commande\_gerant.txt}, support directement lisible ;
  \item les validations mensuelles de quantité et de chiffre d'affaires ;
  \item des graphiques statiques pour le responsable et le contrôle terrain.
\end{itemize}

\section{Principales interfaces graphiques}

\subsection{Accueil et synthèse}

La page d'accueil, présentée à la figure~\ref{fig:ui-accueil}, fournit une lecture
immédiate : quantité à produire, chiffre d'affaires prévu, prochaine fête, part des
prévisions fiables et écart prévu-réel. Elle signale également la fraîcheur des
ventes. Cette alerte est essentielle dans la capture : les données sont arrêtées
au 28 juin 2026, soit quarante jours avant la date d'affichage.

\begin{figure}[H]
\centering
\includegraphics[width=\textwidth]{paul_accueil.png}
\caption{Page d'accueil du tableau de bord PrevisionPaul}
\label{fig:ui-accueil}
\end{figure}

\subsection{Production journalière}

La vue de production du jour affiche les quantités recommandées et permet de
choisir une date, une catégorie et un export Excel. Trois indicateurs résument le
volume du jour, le nombre de références à préparer et l'horizon disponible. La
prévision centrale et la fiabilité restent accessibles dans le tableau afin que
l'utilisateur distingue le besoin moyen de la marge de sécurité.

\begin{figure}[H]
\centering
\includegraphics[width=\textwidth]{paul_production_jour.png}
\caption{Vue de production journalière}
\label{fig:ui-production}
\end{figure}

\subsection{Matières premières}

La vue MRP présente 236 ingrédients distincts pour le mois calculé. Elle regroupe
les solides, les liquides, le budget d'achat estimé, le food-cost plancher et la
couverture des prix. Le tableau détaillé permet de contrôler les matières les plus
importantes. Une seconde section classe les recettes à saisir en priorité selon le
poids matière qu'elles représentent.

\begin{figure}[H]
\centering
\includegraphics[width=\textwidth]{paul_matieres_premieres.png}
\caption{Vue des besoins en matières premières}
\label{fig:ui-mrp}
\end{figure}

\subsection{Événements et commandes}

L'utilisateur peut déclarer un match, un jour férié, un concert ou un autre
événement. Pour un match, l'impact est appris à partir des occurrences passées et
pondéré par l'importance de la rencontre ; pour les autres types, un impact manuel
peut être saisi. L'onglet voisin enregistre les commandes clients datées.

\begin{figure}[H]
\centering
\includegraphics[width=\textwidth]{paul_evenements.png}
\caption{Gestion des événements et des matchs}
\label{fig:ui-evenements}
\end{figure}

\subsection{Assistant déterministe local}

La figure~\ref{fig:ui-assistant} montre l'assistant guidé. L'utilisateur choisit
un type de question et renseigne uniquement les champs nécessaires. La fonction
correspondante lit les fichiers réels et renvoie un texte et, si besoin, un
tableau. Ce choix privilégie la confidentialité et l'absence d'hallucination par
rapport à une conversation libre avec un service externe.

\begin{figure}[H]
\centering
\includegraphics[width=\textwidth]{paul_assistant_local.png}
\caption{Assistant local fondé sur les données calculées}
\label{fig:ui-assistant}
\end{figure}

\subsection{Analyse du chiffre d'affaires et des marges}

L'espace Analyse rassemble l'historique, les prévisions de chiffre d'affaires et
les marges matière par produit. Les cartes indiquent notamment un chiffre
d'affaires historique HT de 51,53 millions MAD sur les données consolidées par le
dashboard, un mois moyen de 873\,399 MAD et les prévisions 2026. Les marges restent
qualifiées d'indicatives tant que toutes les recettes et tous les prix ne sont pas
validés.

\begin{figure}[H]
\centering
\includegraphics[width=\textwidth]{paul_analyse_marges.png}
\caption{Analyse du chiffre d'affaires et des marges matière}
\label{fig:ui-analyse}
\end{figure}

\section{Automatisation et déploiement}

Le dépôt fournit un lanceur double-clic pour l'utilisateur final et un script
d'installation de raccourci avec l'icône PAUL. En production, le code est prévu
dans \texttt{C:\textbackslash PrevisionPaul}, hors des dossiers personnels. Une
tâche Windows nommée « PrevisionPaul - MAJ quotidienne » exécute le script
\texttt{mise\_a\_jour\_quotidienne.py} à 05 h 30. Les étapes et erreurs sont
écrites dans un fichier \texttt{logs/auto\_AAAAMMJJ.log}.

Le code est synchronisé par Git, tandis que les ventes, recettes et prix sensibles
restent exclus du dépôt public. Cette séparation simplifie les mises à jour : un
\texttt{git pull} actualise l'application sans déplacer les données locales.

\section{Tests et validation}

\subsection{Stratégie de test}

La suite Pytest couvre les règles les plus sensibles : préparation des données,
prévisions, fêtes, commandes, incertitude, nomenclatures, coûts, couverture,
assistant et exports. Un test spécifique vérifie que l'assistant ne dépend
d'aucun service réseau. D'autres tests contrôlent l'ajout exact d'une commande,
la neutralisation des pics, les unités de matière et l'absence de régression dans
les tables générées.

La commande \texttt{python -m pytest tests -q} a été exécutée après la
synchronisation GitHub. Les 86 tests ont réussi en 7,85 secondes sur la machine de
validation.

\begin{resultatbox}[title=Résultat de la recette technique,fonttitle=\bfseries]
\centering
\Large\textbf{86 tests réussis sur 86}\par
\normalsize Aucun échec détecté sur le commit \texttt{be291f7}.
\end{resultatbox}

\subsection{Validation mensuelle}

La validation walk-forward disponible rejoue douze mois, de juin 2025 à mai 2026.
Elle confirme que Holt--Winters est le meilleur modèle agrégé parmi les cinq
séries de référence du fichier de validation : WAPE de 8,25\,\% sur les quantités
et 8,28\,\% sur le chiffre d'affaires. Cette mesure ne remplace pas la sélection
par produit ; elle contrôle la cohérence du niveau global.

\begin{table}[H]
\centering
\caption{Résultats de la validation mensuelle agrégée}
\label{tab:validation-mensuelle}
\begin{tabularx}{\textwidth}{Y C{3cm} C{3cm}}
\toprule
\textbf{Modèle} & \textbf{WAPE quantité} & \textbf{WAPE CA} \\
\midrule
Moyenne mobile & 9,36\,\% & 11,09\,\% \\
Tendance & 15,18\,\% & 19,03\,\% \\
Naïf saisonnier & 19,31\,\% & 22,92\,\% \\
Décomposition saisonnière & 16,22\,\% & 20,49\,\% \\
\textbf{Holt--Winters} & \textbf{8,25\,\%} & \textbf{8,28\,\%} \\
\bottomrule
\end{tabularx}
\end{table}

\begin{figure}[H]
\centering
\includegraphics[width=0.92\textwidth]{dashboard_validation_mensuelle.png}
\caption{Graphique de validation mensuelle généré par le pipeline}
\label{fig:validation-mensuelle}
\end{figure}

\subsection{Suivi journalier}

Le suivi du 1\textsuperscript{er} au 28 juin 2026 contient 30\,905 comparaisons
produit-date. À l'échelle de toutes les lignes, le WAPE atteint 43,47\,\%, ce qui
reflète la difficulté de la longue traîne et des petits volumes. Après agrégation
par jour, le WAPE est de 18,62\,\% et le biais de $-18,12\,\%$, indiquant une
sous-prévision globale sur cette fenêtre. Ce diagnostic est affiché plutôt que
masqué : il justifie la marge de sécurité et la nécessité de mettre à jour les
ventes, arrêtées au 28 juin dans la capture du 7 août.

\begin{table}[H]
\centering
\caption{Indicateurs observés dans les exports du 7 juillet 2026}
\label{tab:resultats-observes}
\begin{tabularx}{\textwidth}{p{5.4cm}Y}
\toprule
\textbf{Indicateur} & \textbf{Valeur} \\
\midrule
Produits dans le plan mensuel & 1\,164 \\
Prévision mensuelle centrale totale & 95\,898 unités \\
Quantité mensuelle recommandée totale & 128\,239 unités \\
Part des lignes mensuelles classées fiables & 44,8\,\% \\
Horizon de prévision journalière & 62 jours \\
Produits dans l'export journalier & 1\,109 \\
Ingrédients distincts calculés & 236 \\
Budget matière estimé affiché & 200\,430 MAD \\
Couverture des prix affichée & 94\,\% \\
Tests automatisés & 86 / 86 réussis \\
\bottomrule
\end{tabularx}
\end{table}

\begin{attentionbox}
Les résultats chiffrés correspondent aux fichiers disponibles au 7 juillet 2026.
Ils ne constituent pas une prévision actualisée au 7 août 2026, car les ventes
réelles s'arrêtent au 28 juin. Le dashboard détecte et affiche ce retard. Une
mise à jour SQL ou fichier, suivie d'un recalcul, est nécessaire avant toute
décision opérationnelle réelle.
\end{attentionbox}

\section{Difficultés rencontrées et solutions}

\subsection{Hétérogénéité et qualité des données}

Les exports historiques combinent plusieurs années, sites, encodages et formats.
Les sous-totaux pouvaient doubler les volumes. Nous avons stabilisé un schéma CSV,
introduit des contrôles de seuil, consolidé les mois récents et conservé un
référentiel de nettoyage. Le système refuse désormais de remplacer un historique
valide par une extraction vide.

\subsection{Longue traîne et changements de régime}

Plusieurs centaines de produits ont un historique court ou intermittent. Les
modèles classiques deviennent instables et une erreur relative sur un volume
presque nul perd son sens. Nous avons combiné Croston, des replis simples, un GBM
global et des étiquettes spécifiques. Pour les références qui changent rapidement
de niveau, le modèle journalier augmente le poids de la fenêtre récente tout en
bornant la réaction à un pic isolé.

\subsection{Recettes incomplètes et unités}

La qualité du MRP dépend directement des nomenclatures. Des variations de casse,
des codes internes et des unités incohérentes empêchaient certains rapprochements.
Nous avons ajouté des alias, une normalisation des unités, une gestion récursive
des PSF et un workflow de validation par le chef. Les derniers correctifs du dépôt
traitent notamment une erreur d'échelle kilogramme/gramme et plusieurs familles de
préparations cuisine.

\subsection{Lisibilité de l'interface}

Le dashboard regroupe de nombreux usages dans un seul fichier d'application. Une
navigation plate aurait créé trop d'onglets. Nous avons donc organisé les vues en
cinq groupes, ajouté des sous-onglets thématiques, des aides contextuelles et des
cartes de navigation. Les lectures coûteuses sont mises en cache puis invalidées
après un recalcul.

\section{Conclusion}

La réalisation concrétise l'architecture décrite au chapitre précédent : sources
contrôlées, moteurs séparés, règles métier explicites, exports traçables et
interface orientée décision. Les résultats mensuels sont satisfaisants au niveau
agrégé et la suite automatisée valide les comportements critiques. Le suivi
journalier met toutefois en évidence un biais récent et rappelle qu'un outil de
prévision doit être surveillé et recalibré en continu.
